% Laget automatisk av elevark/lag_elevark.py fra opplegg/newton2-vogn.html. Endringer her blir overskrevet.
\documentclass[11pt]{article}
\usepackage{elevark}
\begin{document}
\elevarkhode{Fysikk 1 · Krefter}{Newtons andre lov med vogn og lodd}{https://fysikklab.no/opplegg/newton2-vogn.html}{Opplegget på nett}

En vogn på en vannrett bane trekkes av et lodd som henger i en snor over en trinse. Hvordan endrer akselerasjonen seg når dere drar hardere, og når vogna blir tyngre? Fra grafene kan dere finne massen til systemet og friksjonen.

\section*{Hensikt}

\begin{itemize}
\item Vise at akselerasjonen er proporsjonal med kraftsummen når massen er fast (del A).
\item Vise at akselerasjonen er omvendt proporsjonal med massen når kraften er fast (del B).
\end{itemize}

\section*{Utstyr}

\begin{itemize}
\item Vogn, vannrett bane med trinse i enden, snor og en loddholder med lodd på 10~g.
\item Vekt, målebånd og stoppeklokke. Fotoporter eller en bevegelsessensor gir mer nøyaktige tider.
\end{itemize}

\section*{Teorien}

Vogna og loddet beveger seg sammen og har samme akselerasjon. Ser vi på dem som ett system med total masse $M$, er den ytre kraften som drar $m_{\mathrm{h}}g$ (tyngden av loddet), og friksjonen $R$ bremser. Newtons andre lov gir $m_{\mathrm{h}}g - R = Ma$. Starter vogna fra ro og kjører strekningen $s$ på tiden $t$, er $a = 2s/t^{2}$.

\section*{Del A: Samme masse, ulik kraft}

\begin{enumerate}
\item Legg alle loddene unntatt ett på vogna, og heng det siste på loddholderen. Vei vogna med alle loddene og loddholderen. Det er den totale massen $M$, og den skal være den samme i hele del A.
\item Mål strekningen $s$ fra startstedet til et merke på banen, for eksempel 80~cm.
\item Slipp vogna fra ro og ta tiden til den passerer merket. Gjør det tre ganger og bruk gjennomsnittet.
\item Flytt ett lodd fra vogna til loddholderen og gjenta. Fortsett til dere har seks målinger.
\end{enumerate}

\section*{Del B: Samme kraft, ulik masse}

\begin{enumerate}
\item Heng et fast lodd på loddholderen, for eksempel 50~g til sammen.
\item Mål tiden over samme strekning $s$ med vogna tom, og så med mer og mer last på vogna.
\item Den totale massen $M$ er vogna, lasten, loddholderen og loddet til sammen.
\end{enumerate}

\section*{Oppgaver}

\begin{oppgaveliste}
\oppgave \fase{gjett}{Gjett}Hva skjer med akselerasjonen hvis dere dobler loddet som henger? Hva skjer hvis dere dobler massen til hele systemet?
\svar
\oppgave \fase{test}{Mål}Gjør del A, eller trykk på «Fyll inn eksempel». Ligger punktene på en rett linje? Hvilken masse gir stigningstallet, og stemmer den med den veide massen?
\svar
\oppgave \fase{forklar}{Forklar}Hvorfor flytter vi loddene fra vogna til loddholderen i stedet for å hente nye lodd? Hvorfor må massen til det hengende loddet regnes med i $M$?
\svar
\oppgave \fase{test}{Regn}Hvor stor er friksjonen ifølge grafen i del A? Hvorfor går ikke linjen gjennom origo?
\svar
\oppgave \fase{test}{Mål}Gjør del B, eller fyll inn eksemplet. Er akselerasjonen omvendt proporsjonal med massen? Sammenlign stigningstallet med tyngden av loddet, og forklar forskjellen.
\svar
\oppgave \fase{utfordring}{Utfordring}Snordraget $S$ er ikke lik tyngden av loddet når systemet akselererer. Bruk Newtons andre lov på loddet alene og regn ut $S$ for den siste målingen i del A. Når blir $S$ nesten lik $m_{\mathrm{h}}g$?
\svar
\end{oppgaveliste}

\end{document}
